\par \textbf{Факторкольца:}
\begin{enumerate}
    \item $\Z / (n)$
    \par $\varphi(x)=x \: mod \: n$ - гомоморфизм из $\Z$ в $\Z_n$. $Im \: \varphi=\Z_n$. $\ker \varphi=n\Z=(n)$. Значит по основной теореме о гомоморфизме $\Z/(n) \cong \Z_n$
    
    \item $K[x]/(x)$
    \par $\varphi(f)=f(0)$ - гомоморфизм из $K[x]$ в $K$, $Im \:  \varphi = K, \ker \varphi=\{x h(x) \: | \: h(x) \in K[x]\}=(x)$. Значит по ОТГ $K[x]/(x) \cong K$
    
    \item $\Z[x] / (n,x)$
    \par Гомоморфизм: $\varphi(f)=f(0) \: mod \: n \Rightarrow \Z[x]/(n,x)=\Z_n$
    
    \item $\Z[x]/(n)$
    \par Гомоморфизм: $\varphi(f)=f \: mod \: n \Rightarrow \Z[x]/(n) \cong \Z_n[x]$
    
    \item $\Z[i]/(1+i)$
    \par $2=(1+i)(1-i) \Rightarrow (2) \subset (1+i) \subset \Z[i]$. Тогда из билета 24 на уд получаем $$(\Z[i]/(2)) / ((1+i) / (2)) \cong \Z_2[i] / (1+i)_{\Z_2} \cong \Z_2$$
    \par Последний переход верен, так как размер $|(1+i)_{\Z_2}|=2$, а $|\Z_2[i]|=4$
    
    \item $\Z[i]/(3+2i)$
    \par Все аналогично предыдущему: $13=(3+2i)(3-2i)$, $|(3+2i)_{\Z_{13}}|=13$ (так как и действительная и мнимая часть взаимно просты с 13), $|\Z_{13}[i]|=13^2 \Rightarrow \Z[i]/(3+2i) \cong \Z_{13}$
\end{enumerate}

\par \textbf{Простые и максимальные идеалы:}
\begin{enumerate}
    \item $I \in \Z$
    \par Так как $\Z$ - КГИ в нем любой простой идеал максимален. Значит $I=(n)$ простой, максимальный тогда и только тогда когда $n$ простое
    
    \item $(1+i)\subset \Z[i]$
    \par $\Z[i]$ евклидово, а значит КГИ (и любой простой идеал максимален). $1+i$ - неразложимый элемент $\Rightarrow$ идеал прост и максимален
    
    \item $(2-5\omega) \subset \Z[\omega]$
    \par Порождающий элемент разложим, а значит идеал непростой и немаксимальный
    
    \item $(n) \subset \Z[w]$
    \par Если $n$ составное, то идеал очевидно не прост и не максимален. В случае простого $n$ фактор по этому идеалу равен $\Z_n$. (???) В $\Z_n$ нет нетривиальных идеалов, а значит, из билета 24 на уд ($K=\Z[\omega], J=(n), I$ - промежуточный идеал) $I/(n)$ либо равно $\Z[\omega]/(n)$, либо равно $(n)/(n)$, а значит $I=(n)$ или $\Z[\omega]$, то есть любой идеал, который содержит $(n)$ - это либо он сам, либо все кольцо, то есть он является простым и максимальным.
\end{enumerate}